\begin{textsample}{POS Dim 1 – ro – Score 55.00 – ro/ro\_ef\_99.txt}  \label{ex:f1_pos_138}
9.9 História – 1º ao 9º Ano 9.9.1 Conceituação do Componente Curricular de História O conhecimento histórico se dá pelo diálogo constante entre o presente e o passado, como destaca a Base Nacional Comum Curricular : todo conhecimento sobre o passado é também um conhecimento do presente elaborado por distintos sujeitos.
O historiador indaga com \textbf{vistas} a identificar, analisar e compreender os significados de diferentes objetos, lugares, circunstâncias, temporalidades, movimentos de pessoas, coisas e saberes.
As perguntas e as elaborações de hipóteses variadas fundam não apenas os marcos de memória, mas também as diversas formas narrativas, ambos expressão do tempo, do caráter social e da prática da produção do conhecimento histórico.
As questões que nos levam a pensar a História como um saber necessário para a formação das crianças, jovens e adultos na escola são as originárias do tempo presente.
O passado, que deve impulsionar a dinâmica do ensino-aprendizagem no Ensino Fundamental, é aquele que dialoga com o tempo \textbf{atual}.
A relação passado/presente não se \textbf{processa} de forma automática, pois \textbf{exige} o conhecimento de referências \textbf{teóricas} capazes de trazer inteligibilidade aos objetos históricos selecionados.
Um objeto só se torna documento quando apropriado por um narrador que a ele confere sentido e significados, tornando -o capaz de expressar a dinâmica da vida das sociedades.
Portanto, o que nos interessa no conhecimento histórico é perceber a forma como os indivíduos construíram, com diferentes linguagens, suas narrações sobre o mundo em que viveram e vivem, suas instituições e organizações sociais.
A Base mostra ainda que a história não \textbf{emerge} como um dado ou um acidente que tudo explica : ela é a \textbf{correlação} de forças, de enfrentamentos e da batalha para a produção de sentidos e significados, que são constantemente reinterpretados por diferentes grupos sociais e suas \textbf{demandas} – o que, consequentemente, suscita outras questões e discussões.
O exercício do “ fazer história ”, de indagar, é marcado, inicialmente, pela constituição de um sujeito.
Em seguida, amplia -se para o conhecimento de um “ Outro ”, às vezes semelhante, muitas vezes diferente.
A percepção do \textbf{aluno} de si e do outro é alargada ainda mais quando se apresentam os conhecimentos sobre outros povos, com seus usos e costumes específicos, ampliando -se para o mundo, sempre em movimento e transformação.
Em meio a \textbf{inúmeras} combinações dessas variáveis – do Eu, do Outro e do Nós –, inseridas em tempos e espaços específicos, indivíduos produzem saberes que os tornam mais aptos para enfrentar situações marcadas pelo conflito ou pela conciliação.
Entre os saberes produzidos, destaca -se a \textbf{capacidade} de comunicação e diálogo, instrumento necessário para o respeito à \textbf{pluralidade} cultural, social e política, bem como para o enfrentamento de circunstâncias marcadas pela tensão e pelo conflito.
A lógica da palavra, da argumentação, é aquela que permite ao sujeito enfrentar os \textbf{problemas} e propor soluções com \textbf{vistas} à superação das contradições políticas, \textbf{econômicas} e sociais do mundo em que \textbf{vivemos}.
Para se pensar o ensino de História, é fundamental considerar a utilização de diferentes fontes e tipos de documentos ( escritos, iconográficos, materiais, imateriais e outros ) capazes de facilitar a \textbf{compreensão} da relação tempo e espaço e das relações sociais que os geraram.
Os registros e vestígios das mais diversas naturezas ( mobiliário, instrumentos de trabalho, música, dentre outros ) deixados pelos indivíduos carregam em si mesmos a experiência humana, as formas específicas de produção, consumo e circulação, \textbf{tanto} de objetos quanto de saberes.
Nessa dimensão, o objeto histórico transforma -se em exercício, em laboratório da memória voltado para a produção de um saber próprio da história.
A utilização de objetos materiais pode \textbf{auxiliar} o professor e os \textbf{alunos} a colocar em questão o significado das coisas do mundo, estimulando a produção do conhecimento histórico em âmbito escolar.
Por meio dessa prática, docentes e \textbf{discentes} poderão \textbf{desempenhar} o papel de agentes do processo de ensino e aprendizagem, assumindo, ambos, uma postura autônoma diante dos conteúdos propostos, no âmbito de um processo adequado ao Ensino Fundamental.
Na organização e desenvolvimento do trabalho pedagógico, o professor mediará a aprendizagem do \textbf{aluno} tendo como referência os processos de construção, a serem contemplados no percurso, devendo garantir : Identificação : Diferentes formas de percepção e interação com um mesmo objeto podem favorecer uma melhor \textbf{compreensão} da história, das mudanças ocorridas no tempo, no espaço e, especialmente, nas relações sociais.
Um mesmo objeto terá diferentes significados, \textbf{dependendo} do tempo, do lugar e da circunstância em que for percebido. \textbf{Contextualização} : Esta é uma \textbf{tarefa} \textbf{imprescindível} para o conhecimento histórico.
Com base em níveis variados de exigência, das operações mais \textbf{simples} às mais elaboradas, os \textbf{alunos} devem ser \textbf{instigados} a aprender a contextualizar.
Saber localizar momentos e lugares específicos de um \textbf{evento}, de um discurso ou de um registro das atividades humanas contribui para evitar atribuição de sentidos e significados não condizentes com uma determinada época, grupo social, comunidade ou território.
Portanto, os estudantes devem identificar, em um contexto, o momento em que uma circunstância histórica é analisada e as condições específicas daquele momento, inserindo o \textbf{evento} em um \textbf{quadro} mais amplo de referências sociais, culturais e \textbf{econômicas}, o que estimula a percepção de que povos e sociedades, em tempos e espaços diferentes, não são tributários dos mesmos valores e princípios da atualidade.
Comparação : Em história, a comparação facilita ver o outro melhor.
Se o tema, por exemplo, for pintura corporal, a comparação entre pinturas de povos indígenas originários e de populações urbanas pode ser bastante esclarecedora quanto ao funcionamento das diferentes sociedades ( as origens das tintas utilizadas, os instrumentos para a realização da pintura e o tempo de duração dos desenhos no corpo... ).
Interpretação : O ensino e aprendizagem em história deve garantir a formação do pensamento \textbf{crítico} do \textbf{aluno}, sendo muito importante para isso que se oportunize o exercício da interpretação, seja de um texto, de um objeto, de uma obra literária, artística ou de um mito, o que \textbf{exige} observação e conhecimento da estrutura do objeto e das suas relações com modelos e formas ( semelhantes ou diferentes ) inseridas no tempo e no espaço.
Interpretações variadas sobre um mesmo objeto tornam mais clara, explícita, a relação sujeito/objeto e, ao mesmo tempo, estimulam a identificação das hipóteses levantadas e dos \textbf{argumentos} selecionados para a comprovação das diferentes proposições.
O exercício da interpretação também permite compreender o significado histórico de uma cronologia e realizar o exercício da composição de outras ordens cronológicas.
Essa prática \textbf{explicita} a \textbf{dialética} da inclusão e da exclusão.
Entre os \textbf{debates} que merecem ser enunciados, destacam -se as dicotomias entre Ocidente e Oriente e os modelos \textbf{baseados} na \textbf{sequência} temporal de \textbf{surgimento}, auge e declínio.
Ambos pretendem dar conta de explicações para questões históricas \textbf{complexas}.
De um \textbf{lado}, a longa existência de tensões ( sociais, culturais, religiosas, políticas e \textbf{econômicas} ) entre sociedades ocidentais e orientais ; de outro, a busca pela \textbf{compreensão} dos modos de organização das várias sociedades que se sucederam ao longo da história. \textbf{Análise} : É uma habilidade bastante \textbf{complexa} porque pressupõe problematizar a própria escrita da história e considerar que, apesar do esforço de organização e de busca de sentido, trata -se de uma atividade em que algo sempre escapa. \textbf{Lembrando} que um dos importantes objetivos de História no Ensino Fundamental é estimular a autonomia de pensamento e a \textbf{capacidade} de reconhecer que os indivíduos agem de acordo com a época e o lugar nos quais vivem, de forma a preservar ou transformar seus hábitos e \textbf{condutas}.
A percepção de que existe uma grande diversidade de sujeitos e histórias estimula o pensamento \textbf{crítico}, a autonomia e a formação para a cidadania.
A busca de autonomia também \textbf{exige} reconhecimento das bases da \textbf{epistemologia} da História, a saber : - a natureza compartilhada do sujeito e do objeto de conhecimento ; - o conceito de tempo histórico em seus diferentes ritmos e durações ; - a concepção de documento como suporte das relações sociais ; - as várias linguagens por meio das quais o ser humano se apropria do mundo.
Sobretudo, é importante destacar que todas as considerações de ordem \textbf{teórica} expostas devem valorizar conhecimentos e experiências dos \textbf{alunos} e professores, tendo em \textbf{vista} a realidade social e o universo da comunidade escolar, bem como seus referenciais históricos, sociais e culturais, principalmente os de significado na sua comunidade.
Na diversidade de \textbf{análises} e proposições, \textbf{espera} -se que os \textbf{alunos} construam as próprias interpretações, de forma fundamentada e rigorosa. \textbf{Convém} destacar as temáticas voltadas para a diversidade cultural e para as múltiplas configurações \textbf{identitárias}, destacando -se as \textbf{abordagens} relacionadas à história dos povos indígenas originários e africanos.
Ressalta -se, também, na formação da sociedade brasileira, a presença de diferentes povos e culturas, suas contradições sociais e culturais e suas articulações com outros povos e sociedades.
A inclusão dos temas obrigatórios definidos pela legislação vigente, tais como a história da África e das culturas afro-brasileira e indígena, deve ultrapassar a dimensão puramente retórica e permitir que se defenda o estudo dessas populações como artífices da própria história do Brasil.
A relevância da história desses grupos humanos reside na possibilidade de os estudantes compreenderem o papel das alteridades presentes na sociedade brasileira, comprometerem -se com elas e, ainda, perceberem que existem outros referenciais de produção, circulação e transmissão de conhecimentos, que podem se entrecruzar com aqueles considerados consagrados nos espaços formais de produção de saber.
Não esquecendo de tratar da percepção estereotipada naturalizada de diferença no tratamento das populações em questão.
A proposta se organiza de forma no sentido do conhecimento histórico ser tratado como uma forma de pensar, entre várias ; uma forma de indagar sobre as coisas do passado e do presente, de construir explicações, \textbf{desvendar} significados, compor e decompor interpretações, em movimento contínuo ao longo do tempo e do espaço.
Enfim, trata -se de transformar a história em ferramenta a serviço de um discernimento maior sobre as experiências humanas e as sociedades em que se vive.
Ao longo da aprendizagem escolar, o \textbf{aluno} vivenciará situações de aprendizagem do conhecimento histórico, em articulação com os conhecimentos dos demais componentes curriculares.
Na prática escolar, o desenvolvimento das \textbf{competências} e habilidades \textbf{esperadas} poderá se dar, favoravelmente, pelas escolhas, por parte do professor, de metodologias ativas, que coloquem o estudante como o \textbf{ator} principal do processo, podendo se dar por meio da : - Busca de informações em diferentes tipos de fontes ( \textbf{entrevistas}, pesquisa bibliográfica, pesquisa na Web, imagens e outras ). - \textbf{Análise} de diferentes tipos de documentos de diferentes naturezas. - Levantamento e troca de informações sobre os objetos de conhecimento. - Comparação de informações e de diferentes interpretações sobre um mesmo acontecimento, fato ou tema histórico. - Formulação de hipóteses e questões a respeito dos temas estudados, como também a busca de possíveis respostas ou soluções. - Trabalho com registros em diferentes formas : textos, livros, fotos, \textbf{vídeos}, \textbf{exposições}, mapas e outros, em diferentes linguagens. - Conhecimento e uso de diferentes medidas de tempo histórico.
Vale destacar a importância da inserção das \textbf{tecnologias} disponíveis, na organização do trabalho pedagógico, notadamente as \textbf{tecnologias} \textbf{digitais}, dado os reconhecidos benefícios da prática da cultura \textbf{digital} na escola, oportunizando ao estudante a integração de conhecimentos por meio da rede, na participação de comunidades de aprendizagem e outras, ampliando suas possibilidades de aprendizagem para além da sala de \textbf{aula}.
A incorporação das \textbf{tecnologias} ao currículo escolar, associada ao conteúdo específico de cada componente, notadamente o de História, favorece o processo de ensino e de aprendizagem no espaço escola e fora dele.
No ambiente escolar, as ações nas quais professores e \textbf{alunos} sejam sujeitos do processo de ensino e aprendizagem devem ser estimuladas.
Nesse sentido, eles próprios devem assumir atitude autônoma diante dos conteúdos propostos no âmbito do Ensino Fundamental.
Considerando esses pressupostos, e em articulação com as \textbf{competências} gerais da BNCC e com as \textbf{competências} específicas da área de Ciências Humanas, o componente curricular de História deve garantir aos \textbf{alunos} o desenvolvimento de \textbf{competências} específicas, a seguir relacionadas.

% matched lemmas: abordagem, aluno, análise, argumento, ator, atual, aula, auxiliar, basear, capacidade, competência, complexo, compreensão, conduta, contextualização, convir, correlação, crítico, debate, demanda, depender, desempenhar, desvendar, dialético, digital, discente, econômico, emergir, entrevista, epistemologia, esperar, evento, exigir, explicitar, exposição, identitária, imprescindível, instigar, inúmero, lado, lembrar, pluralidade, problema, processar, quadro, sequência, simples, surgimento, tanto, tarefa, tecnologia, teórico, vista, vivar, vídeo
\end{textsample}
